% Part: computability
% Chapter: computability-theory
% Section: complete-ce-sets

\documentclass[../../../include/open-logic-section]{subfiles}

\begin{document}

\olfileid{cmp}{thy}{cce}
\olsection{مجموعه‌های محاسبه‌پذیرِ برشمردنیِ کامل}

\begin{defn}
مجموعهٔ $A$ \emph{مجموعه‌ای !!{computably enumerable} و کامل}
(نسبت به کاهش‌پذیری چندبه‌یک) است اگر
\begin{enumerate}
\item $A$ محاسبه‌پذیرِ برشمردنی باشد، و
\item برای هر مجموعهٔ محاسبه‌پذیرِ برشمردنیِ دیگر چون $B$،
  $B\leq_m A$.
\end{enumerate}
\end{defn}

به بیان دیگر، مجموعه‌های محاسبه‌پذیرِ برشمردنیِ کامل «دشوارترین»
مجموعه‌های ممکنِ محاسبه‌پذیرِ برشمردنی‌اند. آن‌ها امکان می‌دهند به
پرسش‌هایی دربارهٔ \emph{هر} مجموعهٔ محاسبه‌پذیرِ برشمردنی پاسخ دهیم.

\begin{thm}
$K$، $K_0$ و~$K_1$ همگی مجموعه‌هایی محاسبه‌پذیرِ برشمردنی و کامل‌اند.
\end{thm}

\begin{proof}
برای دیدن کامل‌بودن $K_0$، بگذارید $B$ هر مجموعهٔ محاسبه‌پذیرِ
برشمردنی باشد. آنگاه برای نمایه‌ای چون $e$،
\[
B = W_e = \Setabs{x}{\cfind{e}(x) \fdefined}.
\]
بگذارید $f$ تابع $f(x)=\tuple{e,x}$ باشد. آنگاه به ازای هر عدد
طبیعی~$x$، $x\in B$ اگر و تنها اگر $f(x)\in K_0$. به بیان دیگر، $f$
مجموعهٔ $B$ را به~$K_0$ کاهش می‌دهد.

برای دیدن کامل‌بودن $K_1$، توجه کنید که در برهان
\olref[k1]{prop:k1} مجموعهٔ $K_0$ را به آن کاهش دادیم. پس بنا بر
\olref[ppr]{prop:trans-red}، هر مجموعهٔ محاسبه‌پذیرِ برشمردنی را نیز
می‌توان به~$K_1$ کاهش داد.

$K_0$ را نیز تقریباً به همان شیوه می‌توان به~$K$ کاهش داد.
\end{proof}

\begin{prob}
کاهشی از $K_0$ به $K$ به دست دهید.
\end{prob}

\begin{digress}
پس معلوم می‌شود همهٔ نمونه‌های مجموعه‌های محاسبه‌پذیرِ برشمردنی که تا
اینجا در نظر گرفته‌ایم یا محاسبه‌پذیرند یا کامل. این باید عجیب به نظر
برسد!{} آیا نمونه‌ای از مجموعهٔ محاسبه‌پذیرِ برشمردنی وجود دارد که نه
محاسبه‌پذیر باشد و نه کامل؟ پاسخ مثبت است، اما این امر تا میانهٔ
دههٔ ۱۹۵۰ اثبات نشد؛ فریدبرگ و موچنیک مستقلاً آن را ثابت کردند.
\end{digress}

\end{document}
